% Generated draft: requires mathematical and source review.
\documentclass[11pt]{article}
\usepackage[margin=25mm]{geometry}
\usepackage{fontspec}
\setmainfont{Latin Modern Roman}
\usepackage{amsmath,amssymb,mathtools,unicode-math}
\setmathfont{Latin Modern Math}
\usepackage{enumitem,needspace,xurl,hyperref}
\setlength{\parindent}{0pt}
\setlength{\parskip}{5pt}
\setlength{\emergencystretch}{3em}
\newcommand{\N}{\mathbb{N}}
\newcommand{\Z}{\mathbb{Z}}
\newcommand{\Q}{\mathbb{Q}}
\newcommand{\R}{\mathbb{R}}
\newcommand{\C}{\mathbb{C}}
\newcommand{\F}{\mathbb{F}}
\newcommand{\A}{\mathbb{A}}
\newcommand{\PP}{\mathbb{P}}
\newcommand{\E}{\mathbb{E}}
\newcommand{\eps}{\varepsilon}
\newcommand{\dd}{\,\mathrm{d}}
\begin{document}
\section*{UnsolvedMath ID: 81}
\textbf{Source problem code:} catalogue entry
\subsection{Definitions and mathematical statement}
This is a research programme rather than a uniquely specified mathematical assertion. A genome space would be a set $G$ of chosen representations of genomes. A metric on $G$ would be a function $d:G\times G\to[0,\infty)$ satisfying, for all $x,y,z\in G$, \[d(x,y)=0\iff x=y,\qquad d(x,y)=d(y,x),\qquad d(x,z)\leq d(x,y)+d(y,z).\] DARPA's challenge is to establish appropriate distance metrics on genome space that incorporate biological utility, rather than merely counting mutations. Neither $G$ nor a mathematical criterion for biological utility or appropriateness is specified in the source. Thus these axioms clarify the word ``metric'' but do not supply the missing biological requirements or turn the agenda into a definite existence theorem. The supplied literature triage dated 2026-08-17 reports active, non-unique work on genome distances and phylogenetic metrics, labels the situation partially solved, and states that theorem-level solved/open classification and a full solution or refutation are not applicable as written. The underlying registry separately labels the record open. Both are supplied, dated catalogue assessments.
\subsection{Short English statement}
Construct biologically meaningful ways to measure how far apart genomes are. A precise problem first requires a genome representation, a biological utility criterion, and the required properties of the distance; these choices are not fixed by the source.
\subsection{Sources}
\begin{itemize}
\item Local UnsolvedMath registry, permanent ID 81, source code catalogue entry, ``The Geometry of Genome Space,'' proposed by DARPA in 2007; supplied statement and background.
\item DARPA DSO, BAA 07-68, Mathematical Challenge Fifteen (2007), the primary source identified by the registry: \url{https://web.math.utk.edu/~vasili/refs/darpa07.MathChallenges.html}. The supplied review describes the evidence as the original one-sentence metric agenda.
\item Supplied literature triage, checked 2026-08-17: classification PARTIAL-PROGRESS; no unique biologically useful genome-space metric specified. It reports no source-text change and cautions that dated triage does not exclude later or unindexed work.
\end{itemize}
\end{document}
